\chapter{2015年度（平成27年度）}
\label{ch:2015}

\section{第1期・専門基礎科目（3題必修）}

\begin{problem}{第1問｜基底と上三角化}{2015-b-1}
$V$ を有限次元実ベクトル空間とする。
\begin{enumerate}
  \item[(1)] 次の定義を述べよ。
  \begin{enumerate}
    \item[(i)] $\{w_1,w_2,\ldots,w_m\}\subset V$ が一次独立であること。
    \item[(ii)] $\{v_1,v_2,\ldots,v_n\}\subset V$ が $V$ の基底であること。
  \end{enumerate}
  \item[(2)] $\{v_1,\ldots,v_n\}$ を $V$ の基底とし，$f:V\to V$ を実ベクトル空間としての同型写像とする。このとき $\{f(v_1),\ldots,f(v_n)\}$ も $V$ の基底であることを示せ。
  \item[(3)] $V=\R^n$ として，$\{v_1,\ldots,v_n\}$ を $V$ の基底とする。各 $i$ $(1\le i\le n)$ に対して $W_i$ を $v_1,\ldots,v_i$ で生成される部分空間とする。$n$ 次正方行列 $A$ が与えられ，各 $i$ に対して $AW_i\subset W_i$ が成り立っているとする。このとき，$n$ 次正則行列 $P$ と $n$ 次上三角行列 $T$ が存在して $P^{-1}AP=T$ とできることを示せ。
\end{enumerate}
\end{problem}
\begin{memo*}
\textbf{解答の見通し。}
一次独立とは、線形結合が0になる方法が係数を全て0にする場合しかないこと、
基底とは、一次独立で空間全体を生成するベクトルの組である。
同型写像は線形な全単射である。線形性で一次結合を $f$ の内側へ入れ、
一次独立性には単射性を、生成することには全射性を使う。

上三角化では、基底ベクトルを列に並べた $P=(v_1\ \cdots\ v_n)$ を作る。
$P^{-1}AP$ は、線形変換 $A$ を基底 $v_1,\ldots,v_n$ の座標で表した行列である。
$Av_i\in\langle v_1,\ldots,v_i\rangle$ を係数表示すると、その係数が
$P^{-1}AP$ の第 $i$ 列になり、第 $i$ 成分より下が全て0になる。
\end{memo*}
\begin{proof}[解答]
\begin{enumerate}
  \item[(1)] (i) $w_1,\ldots,w_m$ が一次独立であるとは
  \[
  \sum_{i=1}^m a_iw_i=0\quad\Longrightarrow\quad a_1=\cdots=a_m=0
  \]
  が任意の $a_i\in\R$ に対して成り立つことをいう。

  (ii) $v_1,\ldots,v_n$ が $V$ の基底であるとは、一次独立であり、かつ
  \[
  V=\langle v_1,\ldots,v_n\rangle
  \]
  となることをいう。

  \item[(2)] $\sum_i a_if(v_i)=0$ なら、線形性と単射性から
  $f(\sum_i a_iv_i)=f(0)$、従って $\sum_i a_iv_i=0$ なので全ての $a_i$ は0。
  また任意の $v\in V$ に対し、全射性から $v=f(w)$ となる $w\in V$ があり、
  $w=\sum_i a_iv_i$ と書けば $v=\sum_i a_if(v_i)$。従って $\{f(v_i)\}$ は基底である。

  \item[(3)] $P=(v_1\ \cdots\ v_n)$ とおく。$\{v_i\}$ は基底なので $P$ は正則である。
  $Av_i\in W_i$ より
  \[
  Av_i=\sum_{j=1}^i t_{ji}v_j.
  \]
  $T=(t_{ji})$ とおけば $T$ は上三角であり、
  \[
  AP=(Av_1\ \cdots\ Av_n)=PT.
  \]
  従って $P^{-1}AP=T$。
\end{enumerate}
\end{proof}

\begin{problem}{第2問｜広義積分}{2015-b-2}
$a$ を実数とする。区間 $(0,\infty)$ 上の関数 $f_a(x)$ を
\[
f_a(x)=x^a\log x
\]
によって定める。
\begin{enumerate}
  \item[(1)] 関数 $f_a(x)$ の不定積分を求めよ。
  \item[(2)] 正数 $R>0$ に対して
  \[
  I(a,R)=\iint_{B(R)}f_a(1+x^2+y^2)\,\d x\d y
  \]
  とおく。ただし $B(R)$ は $xy$ 平面上の原点を中心とする半径 $R$ の円の内部とする。このとき $I(0,1)$ の値を求めよ。
  \item[(3)] (2) で定めた $I(a,R)$ に対して $\displaystyle\lim_{R\to\infty}I(a,R)$ が収束するような $a$ の範囲を求めよ。
\end{enumerate}
\end{problem}
\begin{memo*}
\textbf{解答の見通し。}
不定積分では $\log x$ を微分すると $1/x$ になり、$x^a$ は積分しやすいので部分積分を使う。
$\int x^a\,\d x=x^{a+1}/(a+1)$ には $a\ne-1$ が必要なため、$a=-1$ だけ別に計算する。

二重積分の領域は原点中心の円板で、被積分関数も $x^2+y^2$ だけで決まる。
そこで極座標を使うと面積要素は $\d x\d y=r\,\d r\d\theta$ となり、
角度積分は係数 $2\pi$ になる。$u=1+r^2$ と置けば
不定積分と同じ $\int u^a\log u\,\d u$ へ戻る。
最後は $t=\log u$ と置換する。$\log$ が $t$ に化けて $\int_0^\infty te^{(a+1)t}\,\d t$ になるので、
指数が減衰するかどうか、すなわち $a+1$ の符号だけを見ればよい。
\end{memo*}
\begin{proof}[解答]
\begin{enumerate}
  \item[(1)]
  \[
  \int x^a\log x\,\d x=
  \begin{cases}
  \displaystyle
  x^{a+1}\left(\frac{\log x}{a+1}-\frac1{(a+1)^2}\right)+C &(a\ne-1),\\[2mm]
  \displaystyle\frac12(\log x)^2+C &(a=-1).
  \end{cases}
  \]

  \item[(2)] 極座標と $u=1+r^2$ により
  \[
  I(a,R)=2\pi\int_0^R(1+r^2)^a\log(1+r^2)r\,\d r
  =\pi\int_1^{1+R^2}u^a\log u\,\d u.
  \]
  従って
  \[
  I(0,1)=\pi[u\log u-u]_1^2=\pi(2\log2-1).
  \]

  \item[(3)] $t=\log u$ とおくと
  \[
  \int_1^\infty u^a\log u\,\d u
  =\int_0^\infty te^{(a+1)t}\,\d t.
  \]
  これは $a+1<0$ のとき、かつそのときに限り収束する。従って
  \[
  a<-1.
  \]
  なお、このとき極限値は $\pi/(a+1)^2$ である。
\end{enumerate}
\end{proof}

\begin{problem}{第3問｜二重周期関数}{2015-b-3}
連続写像 $f:\R^2\to\R$ が，任意の $(x,y)\in\R^2$ について
$f(x+1,y)=f(x,y+1)=f(x,y)$ を満たすとする。このとき，以下の問いに答えよ。
\begin{enumerate}
  \item[(1)] $f(\R^2)=f([0,1]\times[0,1])$ を示せ。
  \item[(2)] 写像 $f$ は最大値と最小値を持つことを示せ。
  \item[(3)] 写像 $f$ の最大値を $b$，最小値を $a$ とする。このとき，連続写像 $\sigma:[0,1]\to\R^2$ で，合成 $f\circ\sigma:[0,1]\to[a,b]$ が全射となるようなものが存在することを示せ。
\end{enumerate}
\end{problem}
\begin{memo*}
\textbf{解答の見通し。}
$\lfloor x\rfloor$ は $x$ 以下の最大の整数であり、$0\le x-\lfloor x\rfloor<1$ を満たす。
各座標から整数部分を引けば、周期性により関数値を変えずに任意の点を
基本領域 $[0,1]^2$ へ移せる。
$[0,1]^2$ は閉かつ有界なのでコンパクトであり、連続関数はそこで最大値・最小値を取る。

最小値を取る点 $p$ と最大値を取る点 $q$ を線分 $c(t)=(1-t)p+tq$ で結ぶ。
この線分は正方形の中に残り、一変数連続関数 $f(c(t))$ は端点で最小値と最大値を取る。
従って中間値の定理により、その間の全ての値を通る。
\end{memo*}
\begin{proof}[解答]
\begin{enumerate}
  \item[(1)] $m=\lfloor x\rfloor$, $n=\lfloor y\rfloor$ とすると、周期性より
  \[
  f(x,y)=f(x-m,y-n),\qquad (x-m,y-n)\in[0,1)^2.
  \]
  従って $f(\R^2)\subset f([0,1]^2)$。逆の包含は明らかなので等号が成り立つ。

  \item[(2)] $[0,1]^2$ はコンパクトだから、$f$ はそこで最大値と最小値を取る。
  (1) より、それらは $\R^2$ 全体での最大値・最小値でもある。

  \item[(3)] $f(p)=a$, $f(q)=b$ となる $p,q\in[0,1]^2$ を取り
  \[
  \sigma(t)=(1-t)p+tq\qquad(0\le t\le1)
  \]
  とおく。$f\circ\sigma$ は連続で端点の値が $a,b$ だから、
  中間値の定理より $[a,b]$ への全射である。
\end{enumerate}
\end{proof}

\section{専門科目（8題から2題選択）}

\begin{problem}{第1問｜有限体上の行列群}{2015-s-1}
$\mathbb F$ を素数 $p$ を位数とする有限体とする。$GL_n(\mathbb F)$ を $\mathbb F$ の元を成分とする
$n$ 次正則行列全体とし，
\[
G=\{A=(a_{ij})\in GL_3(\mathbb F)\mid\det A=1,\ a_{12}=0,\ a_{22}=1,\ a_{32}=0\}
\]
とする。
\begin{enumerate}
  \item[(1)] $G$ が行列の積に関して群をなすことを示せ。
  \item[(2)] 写像 $(a_{ij})\mapsto\begin{pmatrix}a_{11}&a_{13}\\a_{31}&a_{33}\end{pmatrix}$ により群準同型 $\varphi:G\to GL_2(\mathbb F)$ が得られることを示せ。
  \item[(3)] (2) の写像 $\varphi$ の核を $K$ とする。$G$ における $K$ の中心化群 $C_G(K)$ を求めよ。
  \item[(4)] $X=\begin{pmatrix}1&0&0\\0&1&0\\1&0&1\end{pmatrix}$ とし，$K$ と $X$ で生成される部分群を $H$ とする。$H$ の位数を求めよ。また，$H$ の中心 $Z(H)$ を求めよ。
\end{enumerate}
\end{problem}
\begin{memo*}
\textbf{解答の見通し。}
基底順を $(e_1,e_3;e_2)$ に並べ替えると、元は
線形部分 $B\in SL_2(\mathbb F_p)$ と平行移動部分 $u\in\mathbb F_p^2$ の組 $(B,u)$ になる。
この積公式で群・準同型・核をまとめて扱う。
中心化群では全ての $u$ を固定する $B$ を求め、
$H$ は $K$ と位数 $p$ の元による半直積として数える。
\end{memo*}
\begin{proof}[解答]
(1) 基底順を $(e_1,e_3;e_2)$ と見れば元は
$\begin{pmatrix}B&0\\u&1\end{pmatrix}$（$B\in SL_2,u\in\mathbb F_p^2$）と書ける。
この表示で積と逆元は
\[
\begin{pmatrix}B&0\\u&1\end{pmatrix}
\begin{pmatrix}C&0\\v&1\end{pmatrix}
=\begin{pmatrix}BC&0\\uC+v&1\end{pmatrix},
\]
\[
\begin{pmatrix}B&0\\u&1\end{pmatrix}^{-1}
=\begin{pmatrix}B^{-1}&0\\-uB^{-1}&1\end{pmatrix}.
\]
$BC,B^{-1}\in SL_2(\mathbb F_p)$ なので $G$ は群である。

(2) $\varphi(B,u)=B$ について $\varphi((B,u)(C,v))=BC=\varphi(B,u)\varphi(C,v)$ だから、$\varphi$ は群準同型である。

(3) $K=\ker\varphi=\{(I,u)\}\cong\mathbb F_p^2$ である。$(B,w)$ が全ての $(I,u)\in K$ と可換する条件を積の公式へ代入すると
\[
(B,w)(I,u)=(B,w+u),\qquad
(I,u)(B,w)=(B,uB+w).
\]
従って全ての行ベクトル $u$ について $uB=u$、すなわち $B=I$ が必要十分である。
よって $C_G(K)=K$。

(4) $B_0=\varphi(X)=\begin{pmatrix}1&0\\1&1\end{pmatrix}$ とおくと
\[
B_0^j=\begin{pmatrix}1&0\\j&1\end{pmatrix}\qquad(j\in\mathbb F_p)
\]
なので $B_0$ の位数は $p$。$K\cap\langle X\rangle=\{1\}$ で、$X$ は共役により
$K$ を保つから、$H=K\rtimes\langle X\rangle$ であり
$|H|=p^2p=p^3$ となる。

中心の元は特に $K$ と可換するので、上の議論から $(I,u)$ の形に限る。
$u=(c,d)$ とすると
\[
(I,u)(B_0,0)=(B_0,uB_0),\qquad
(B_0,0)(I,u)=(B_0,u),
\]
かつ $uB_0=(c+d,d)$ だから、$X$ と可換する条件は $d=0$ である。従って中心は
\[
Z(H)=\left\{\begin{pmatrix}1&0&0\\c&1&0\\0&0&1\end{pmatrix}\middle|c\in\mathbb F_p\right\},
\]
従って位数 $p$。
\end{proof}

\begin{problem}{第2問｜整数係数多項式環}{2015-s-2}
整数係数1変数多項式環 $\Z[X]$ から実数環の直積環 $\R\times\R$ への環準同型写像 $\Phi$ を，
$f(X)\in\Z[X]$ に対し
\[
\Phi(f(X))=\bigl(f(\sqrt7),\,f(\sqrt{13})\bigr)
\]
で定める。
\begin{enumerate}
  \item[(1)] $f(X)\in\Z[X]$ に対し，$f(\sqrt7)=0$ となるなら，$f(X)$ は $X^2-7$ で割り切れることを示せ。
  \item[(2)] $\ker\Phi=((X^2-7)(X^2-13))$ を示せ。
  \item[(3)] $\Phi(X^2-7)$ で生成される $\operatorname{Im}\Phi$ のイデアルを $P_1$，$\Phi(X^2-13)$ で生成される $\operatorname{Im}\Phi$ のイデアルを $P_2$ とする。$P_1$ および $P_2$ は $\operatorname{Im}\Phi$ の素イデアルであるが極大イデアルではないことを示せ。
  \item[(4)] $\operatorname{Im}\Phi$ の素イデアルで，$P_1$ と $P_2$ の両方を含むものの個数を求めよ。
\end{enumerate}
\end{problem}
\begin{memo*}
\textbf{解答の見通し。}
$X^2-7$ で割った余りは一次式なので、$\sqrt7$ の無理性から余りが0と分かる。
この議論を $\sqrt{13}$ にも順に使って核を決める。
$P_1,P_2$ の素性は各成分への射影の商が整域になることから示す。
両方を含む素イデアルは差6も含むため、標数2と3に分けて因数分解を数える。
\end{memo*}
\begin{proof}[解答]
(1) 多項式の除法により
\[
f(X)=(X^2-7)q(X)+rX+s\qquad(r,s\in\Z)
\]
と書ける。$f(\sqrt7)=0$ なら $r\sqrt7+s=0$。もし $r\ne0$ なら
$\sqrt7=-s/r\in\Q$ となって矛盾するので $r=0$、さらに $s=0$ である。
従って $X^2-7$ は $f$ を割る。

(2) $f\in\ker\Phi$ とする。(1) から $f=(X^2-7)q$ と書ける。
$0=f(\sqrt{13})=6q(\sqrt{13})$ なので $q(\sqrt{13})=0$。
同じ議論により $X^2-13$ は $q$ を割るから
\[
f\in((X^2-7)(X^2-13)).
\]
逆にこの積の倍数は $\sqrt7$, $\sqrt{13}$ の両方で0になるので、
\[
\ker\Phi=((X^2-7)(X^2-13)).
\]

(3) $\operatorname{Im}\Phi$ から第1成分を取る写像の像は $\Z[\sqrt7]$ で、核は
$P_1$ である。従って準同型定理から
\[
\operatorname{Im}\Phi/P_1\cong\Z[\sqrt7].
\]
これは整域なので $P_1$ は素イデアルだが、例えば2は逆元を持たないので体ではなく、
$P_1$ は極大でない。第2成分への射影
$\Z[\sqrt7]\times\Z[\sqrt{13}]\to\Z[\sqrt{13}]$
を $\Phi$ と合成した写像の核は $P_2$、像は $\Z[\sqrt{13}]$ である。従って
$\operatorname{Im}\Phi/P_2\cong\Z[\sqrt{13}]$ であり、$P_2$ も素だが極大でない。

(4) $P_1,P_2$ の両方を含む素イデアル $\mathfrak p$ は
\[
\Phi(X^2-7)-\Phi(X^2-13)=\Phi(6)
\]
を含む。従って $6\in\mathfrak p$ であり、素性から $2\in\mathfrak p$ または
$3\in\mathfrak p$ である。法2では二つの関係式はいずれも
\[
X^2-1=(X-1)^2
\]
となるので素イデアルは $(2,X-1)$ に対応する1個。法3では
\[
X^2-1=(X-1)(X+1)
\]
となるので $(3,X-1)$ と $(3,X+1)$ に対応する2個である。
従って求める素イデアルは合計 $1+2=3$ 個である。
\end{proof}

\begin{problem}{第3問｜リフトと基本群}{2015-s-3}
$X,Y$ を弧状連結な位相空間，$x_0$ を $X$ の点とする。連続写像 $f:Y\to X$ は次の性質を満たすとする。
\begin{itemize}
  \item $F=f^{-1}(\{x_0\})$ とおく。$F$ の相対位相は離散位相である。
  \item $A$ を $[0,1]$（または $[0,1]\times[0,1]$）とする。このとき $g:A\to X$ と $f(b)=g(0)$（または $g(0,0)$）となる $b\in Y$ に対して，連続写像 $\tilde g:A\to Y$ で $\tilde g(0)=b$（または $\tilde g(0,0)=b$）かつ $f\circ\tilde g=g$ となるものが唯一つ存在する。この $\tilde g$ を $g$ のリフトと呼ぶ。
\end{itemize}
以下の問いに答えよ。
\begin{enumerate}
  \item[(1)] $f_*:\pi_1(Y,y_0)\to\pi_1(X,x_0)$ は単射であることを示せ。ただし，$y_0\in F$ であり，$\pi_1(X,x_0),\pi_1(Y,y_0)$ はそれぞれの位相空間の基本群（1次ホモトピー群）であり，$f_*:\pi_1(Y,y_0)\to\pi_1(X,x_0)$ は $f_*([\delta])=[f\circ\delta]$ によって定められる基本群の間の準同型写像である。
  \item[(2)] $\varphi:\pi_1(X,x_0)\to F$ を $\varphi([\gamma])=\tilde\gamma(1)$ で定めるとき，以下に答えよ。ただし $\tilde\gamma$ は $\tilde\gamma(0)=y_0$ を満たす $\gamma$ のリフトである。\\
  (i) $\varphi$ の定義はホモトピー類の代表元の取り方によらないことを示せ。\\
  (ii) $\varphi$ は全射であることを示せ。\\
  (iii) $\varphi([\gamma_0])=\varphi([\gamma_1])$ が成り立つことと $[\gamma_0][\gamma_1]^{-1}\in f_*(\pi_1(Y,y_0))$ は同値であることを示せ。
\end{enumerate}
\end{problem}
\begin{memo*}
\textbf{解答の見通し。}
ホモトピーを持ち上げると、その端点はファイバー $F$ の中を連続に動く。
$F$ は離散、パラメータ区間は連結なので端点は動けない、という事実を繰り返し使う。
$f_*$ の単射性、$\varphi$ の代表元によらないこと、二つのリフトの終点条件を
全てこの一つの原理から示す。
\end{memo*}
\begin{proof}[解答]
(1) $y_0$ を基点とする閉道 $\eta:[0,1]\to Y$ に対して
\[
f_*([\eta])=[f\circ\eta]=1
\]
と仮定する。このとき $f\circ\eta$ と定値道 $c_{x_0}$ の間に、端点を固定するホモトピー
\[
H:[0,1]\times[0,1]\to X
\]
が存在する。$H(0,t)=f(\eta(t))$ とし、このホモトピーを初期リフト $\eta$ から持ち上げて
$\widetilde H:[0,1]^2\to Y$ を得る。

各 $s\in[0,1]$ について $H(s,0)=H(s,1)=x_0$ だから
\[
\widetilde H(s,0),\ \widetilde H(s,1)\in F=f^{-1}(x_0).
\]
二つの写像 $s\mapsto\widetilde H(s,0)$ と $s\mapsto\widetilde H(s,1)$ は連続であり、
$[0,1]$ は連結、$F$ は離散なので、それぞれ一定である。$s=0$ では両端点が $y_0$ だから
\[
\widetilde H(s,0)=\widetilde H(s,1)=y_0
\qquad(0\le s\le1).
\]
さらに $H(1,t)=x_0$ であり、その $y_0$ から始まるリフトは道のリフトの一意性により
定値道 $c_{y_0}$ である。従って $\eta=\widetilde H(0,\cdot)$ は
$c_{y_0}=\widetilde H(1,\cdot)$ と端点を固定してホモトピックであり、$[\eta]=1$ である。
よって $\ker f_*=\{1\}$、すなわち $f_*$ は単射である。

(2)(i) $x_0$ を基点とする閉道 $\gamma$ の $y_0$ から始まる一意なリフトを
$\widetilde\gamma$ とし
\[
\varphi([\gamma])=\widetilde\gamma(1)\in F
\]
と定める。同じ類を表す閉道 $\gamma_0,\gamma_1$ の間の端点固定ホモトピーを持ち上げると、
リフトの終点は $F$ の中を連続に動く。$F$ は離散なので終点は一定であり
\[
\widetilde\gamma_0(1)=\widetilde\gamma_1(1).
\]
従って $\varphi$ は代表元の取り方によらず well-defined である。

(ii) 任意の $y\in F$ に対して、$Y$ の弧状連結性より
$y_0$ から $y$ への道 $\alpha$ を取れる。このとき
\[
\gamma=f\circ\alpha
\]
は $x_0$ を基点とする閉道であり、$\alpha$ 自身が $\gamma$ の $y_0$ から始まるリフトである。
従って $\varphi([\gamma])=\alpha(1)=y$ となり、$\varphi$ は全射である。

(iii) $\gamma_0,\gamma_1$ の $y_0$ から始まるリフトを
$\widetilde\gamma_0,\widetilde\gamma_1$ とする。もし
\[
\varphi([\gamma_0])=\varphi([\gamma_1])=:y
\]
なら、道
\[
\widetilde\gamma_0*\widetilde\gamma_1^{-1}
\]
は $y_0$ を基点とする $Y$ の閉道であり、その $f$ による像は
$\gamma_0*\gamma_1^{-1}$ である。従って
\[
[\gamma_0][\gamma_1]^{-1}
=f_*([\widetilde\gamma_0*\widetilde\gamma_1^{-1}])
\in f_*(\pi_1(Y,y_0)).
\]

逆に $[\gamma_0][\gamma_1]^{-1}\in f_*(\pi_1(Y,y_0))$ とする。
このとき $\gamma_0*\gamma_1^{-1}$ は、ある $Y$ の閉道の像と端点を固定してホモトピックである。
ホモトピーを持ち上げ、上と同じく終点が離散集合 $F$ の中で一定であることを使うと、
$\gamma_0*\gamma_1^{-1}$ の $y_0$ から始まるリフトも $y_0$ で終わる。
このリフトの後半を逆向きに見ると、それは $\gamma_1$ の $y_0$ から始まるリフトであり、
その終点は $\widetilde\gamma_0(1)$ である。道のリフトの一意性より、このリフトは
$\widetilde\gamma_1$ と一致する。従って
\[
\widetilde\gamma_0(1)=\widetilde\gamma_1(1),
\]
すなわち $\varphi([\gamma_0])=\varphi([\gamma_1])$ である。
\end{proof}

\begin{problem}{第4問｜楕円放物面}{2015-s-4}
$\R^2$ 上の関数 $f$ を $f(x,y)=2x^2+y^2$ で定める。$f$ のグラフを $S$ で表し，
$S$ 内の曲線 $C$ を $C=\{(t,t,f(t,t))\mid t\in\R\}$ で定義する。
\begin{enumerate}
  \item[(1)] $S$ の各点でのガウス曲率および平均曲率を求めよ。
  \item[(2)] $C$ の接線に関する $S$ の法曲率を求めよ。
  \item[(3)] (2) で求めた法曲率が $S$ の平均曲率に等しいような $C$ の点を全て求めよ。
\end{enumerate}
\end{problem}
\begin{memo*}
\textbf{解答の見通し。}
グラフ表示から第1・第2基本形式の係数を順に計算する。
Gauss 曲率と平均曲率はこれらの標準公式へ代入する。
曲線 $C(t)=X(t,t)$ の座標速度は $(1,1)$ なので、
法曲率は第2基本形式と第1基本形式の比で求まる。
最後は $x=y=t$ を代入し、分母が正であることを確認して方程式を整理する。
\end{memo*}
\begin{proof}[解答]
(1) グラフを $X(x,y)=(x,y,2x^2+y^2)$ とパラメータ表示する。
\[
X_x=(1,0,4x),\qquad X_y=(0,1,2y)
\]
なので、第1基本形式の係数と上向き単位法線は
\[
E=1+16x^2,\quad F=8xy,\quad G=1+4y^2,
\]
\[
n=\frac{(-4x,-2y,1)}W,\qquad W=\sqrt{1+16x^2+4y^2}.
\]
また $X_{xx}=(0,0,4)$, $X_{xy}=0$, $X_{yy}=(0,0,2)$ より
第2基本形式の係数は
\[
e=\frac4W,\qquad f=0,\qquad g=\frac2W.
\]
従って
\begin{align*}
K&=\frac{eg-f^2}{EG-F^2}
=\frac{8/W^2}{W^2}=\frac8{W^4},\\
H&=\frac{eG-2fF+gE}{2(EG-F^2)}\\
&=\frac{4(1+4y^2)+2(1+16x^2)}{2W^3}
=\frac{3+16x^2+8y^2}{W^3}.
\end{align*}
すなわち
\[
K=\frac8{W^4},\qquad H=\frac{3+16x^2+8y^2}{W^3}.
\]
(2) $C(t)=X(t,t)$ の座標方向は $(\d x,\d y)=(1,1)$ なので、法曲率は
第2基本形式と第1基本形式の比から
\begin{align*}
k_n(t)
&=\frac{e+2f+g}{E+2F+G}\bigg|_{(x,y)=(t,t)}\\
&=\frac{6/\sqrt{1+20t^2}}{2+36t^2}\\
&=\frac6{\sqrt{1+20t^2}(2+36t^2)}.
\end{align*}
従って
\[
k_n(t)=\frac6{\sqrt{1+20t^2}(2+36t^2)}.
\]
(3) 平均曲率に $x=y=t$ を代入すると
\[
H(t,t)=\frac{3+24t^2}{(1+20t^2)^{3/2}},\qquad
k_n(t)=\frac6{\sqrt{1+20t^2}(2+36t^2)}.
\]
両辺に正の $(1+20t^2)^{3/2}(2+36t^2)$ を掛けて整理すると
\begin{align*}
6(1+20t^2)&=(3+24t^2)(2+36t^2)\\
&=6+156t^2+864t^4,
\end{align*}
すなわち $36t^2(1+24t^2)=0$。実数では $t=0$ だけなので、点は $(0,0,0)$ のみ。
\end{proof}

\begin{problem}{第5問｜3次方程式}{2015-s-5}
複素変数の3次関数 $f(z)=z^3+az+b$ について，次の問いに答えよ。
\begin{enumerate}
  \item[(1)] $R:=\max\{|a|,|b|\}+1$ とおく。このとき，$|z|\ge R+1$ を満たす $z$ に対して，$f(z)\ne0$ を示せ。
  \item[(2)] 
  \[
  P_n:=P_n(a,b)=\frac1{2\pi i}\int_{|z|=R+1}\frac{z^n(3z^2+a)}{z^3+az+b}\,\d z
  \qquad(n\ge0)
  \]
  とおく。このとき，$\Delta:=P_0P_3^2-\dfrac12P_2^3$ を $a,b$ を用いて表せ。
  \item[(3)] $\Delta=0$ かつ $(a,b)\ne(0,0)$ のとき，3次方程式 $z^3+az+b=0$ の解を全て求めよ。
\end{enumerate}
\end{problem}
\begin{memo*}
\textbf{解答の見通し。}
大円上では最高次項 $z^3$ が $az+b$ より大きいことを示し、全ての根を円内に入れる。
$f'/f$ の積分は偏角の原理により根の冪和 $P_n=\sum\alpha_j^n$ になる。
Vieta の公式で $P_2,P_3$ を $a,b$ に直し、
$\Delta=0$ では重根候補 $r=-3b/(2a)$ を作って多項式を因数分解する。
\end{memo*}
\begin{proof}[解答]
(1) $M=\max\{|a|,|b|\}$ とおくと $R=M+1$ である。$t=|z|\ge R+1=M+2$ なら
\[
|az+b|\le |a|t+|b|\le M(t+1)\le(t-2)(t+1)<t^3=|z^3|.
\]
従って逆三角不等式から
$|f(z)|\ge |z|^3-|az+b|>0$ となり、$f(z)\ne0$ である。

(2) $f'(z)=3z^2+a$ なので、偏角原理（対数微分の留数）から
\[
P_n=\frac1{2\pi i}\int_Cz^n\frac{f'(z)}{f(z)}\,\d z
=\sum_{j=1}^3\alpha_j^n,
\]
ここで $\alpha_1,\alpha_2,\alpha_3$ は重複度を込めた3根である。
$P_0=3$ であり、$\alpha_1+\alpha_2+\alpha_3=0$、
$\alpha_1\alpha_2+\alpha_2\alpha_3+\alpha_3\alpha_1=a$、
$\alpha_1\alpha_2\alpha_3=-b$ を使うと
\begin{align*}
P_2&=(\alpha_1+\alpha_2+\alpha_3)^2
-2(\alpha_1\alpha_2+\alpha_2\alpha_3+\alpha_3\alpha_1)=-2a,\\
P_3&=(\alpha_1+\alpha_2+\alpha_3)^3
-3(\alpha_1+\alpha_2+\alpha_3)
(\alpha_1\alpha_2+\alpha_2\alpha_3+\alpha_3\alpha_1)
+3\alpha_1\alpha_2\alpha_3=-3b.
\end{align*}
従って
\[
\Delta=3(-3b)^2-\frac12(-2a)^3=4a^3+27b^2.
\]

(3) $\Delta=0$ かつ $(a,b)\ne(0,0)$ なら $a\ne0$ である。
$r=-3b/(2a)$ とおく。$4a^3+27b^2=0$ を用いると
\[
-3r^2=a,\qquad 2r^3=b.
\]
従って係数を比較して
\[
(z-r)^2(z+2r)=z^3-3r^2z+2r^3=z^3+az+b=f(z).
\]
よって根は $r=-3b/(2a)$（重複度2）と $-2r=3b/a$（重複度1）である。
\end{proof}

\begin{problem}{第6問｜常微分方程式}{2015-s-6}
\begin{enumerate}
  \item[(1)] 微分方程式の初期値問題
  \[
  \begin{cases}
  \dfrac{\d y}{\d x}=\sin y,\\[1mm]
  y(0)=\dfrac\pi2
  \end{cases}
  \]
  の解を求めよ。
  \item[(2)] $f(x,y)$ を $\R^2$ 上の $C^1$ 級関数とする。微分方程式
  \[
  \frac{\d y}{\d x}=\sin f(x,y)
  \]
  の任意の解の定義域は $\R$ 全体となることを示せ。
\end{enumerate}
\end{problem}
\begin{memo*}
\textbf{解答の見通し。}
初めの方程式は変数分離し、$1/\sin y$ の原始関数を
$\log\tan(y/2)$ として求める。
一般の方程式では右辺の絶対値が1以下なので、解は有限時間に高々その時間幅しか動かない。
従って有限端点で発散できず、局所存在・一意性定理によりさらに延長できる。
\end{memo*}
\begin{proof}[解答]
(1) $0<y<\pi$ の範囲で変数分離すると
\[
\frac{\d y}{\sin y}=\d x,\qquad
\int\frac{\d y}{\sin y}=\log\tan\frac y2.
\]
ここで原始関数は
\[
\frac{\d}{\d y}\log\tan\frac y2
=\frac{1}{\tan(y/2)}\cdot\frac12\sec^2\frac y2
=\frac1{2\sin(y/2)\cos(y/2)}
=\frac1{\sin y}
\]
から確かめられる。
従って $\log\tan(y/2)=x+C$。初期条件 $y(0)=\pi/2$ から $C=0$ なので
\[
\tan\frac y2=e^x,\qquad y(x)=2\arctan(e^x).
\]
(2) 以下、解とは延長できないもの、すなわち極大解を指す。
$y'=\sin f(x,y)$ では $|y'|\le1$ である。解が $x_0$ から有限時刻 $b$ まで存在するとき
\[
|y(x)-y(x_0)|\le|x-x_0|
\]
だから $y$ は有限区間で発散しない。右辺は $C^1$ なので局所存在・一意性定理により
有限な端点の近くからさらに解を延長できる。従って極大解の定義域の端点は有限にならず、定義域は $\R$ 全体である。
\end{proof}

\begin{problem}{第7問｜収束定理}{2015-s-7}
$p>0$ を定数とする。測度空間 $(X,\mathfrak B,\mu)$ 上の $p$ 乗可積分関数 $f$ は
$0<f<\infty$ かつ $\int_Xf^p\,\d\mu=1$ を満たしているとする。このとき次の値を求めよ。
\[
\lim_{n\to\infty}\int_X n\left|\sin\left(\frac{f(x)}n\right)\right|^p\,\d\mu(x).
\]
（収束定理を使う場合はその根拠を述べよ。）
\end{problem}
\begin{memo*}
\textbf{解答の見通し。}
被積分関数の $n$ の効き方を先に取り出す。$u_n=f(x)/n$ とおくと
$|\sin u_n|^p=u_n^p(|\sin u_n|/u_n)^p$ であり、$u_n^p=f(x)^p/n^p$ だから
被積分関数は $n^{1-p}f(x)^p$ に、1へ収束する因子を掛けた形になる。
従って $n^{1-p}$ の挙動、すなわち $p$ と1の大小で結論が三つに分かれる。

$|\sin u|\le|u|$ による上からの評価 $n|\sin u_n|^p\le n^{1-p}f^p$ は、
$p>1$ ならそのまま積分して0を与える。$p=1$ ではこの評価が可積分な優関数 $f$ を与えるので
優収束定理が使える。$p<1$ では被積分関数が各点で $\infty$ へ発散するので、
下から評価する Fatou の補題を使う。
\end{memo*}
\begin{proof}[解答]
各 $x\in X$ について $u_n=f(x)/n$ とおく。$0<f(x)<\infty$ だから $u_n>0$ かつ $u_n\to0$ であり、
$\sin u/u\to1$ $(u\to0)$ より
\[
\frac{|\sin u_n|}{u_n}\longrightarrow1.
\]
$|\sin u_n|^p=u_n^p\left(|\sin u_n|/u_n\right)^p$ と $u_n^p=f(x)^p/n^p$ を使うと
\[
n\left|\sin\left(\frac{f(x)}n\right)\right|^p
=n^{1-p}f(x)^p
\left(\frac{|\sin u_n|}{u_n}\right)^p
\tag{$*$}
\]
であり、右端の因子は各点で1へ収束する。
また $u\ge0$ で $|\sin u|\le u$ だから
\[
0\le n\left|\sin\left(\frac{f(x)}n\right)\right|^p
\le n\,u_n^p=n^{1-p}f(x)^p
\tag{$**$}
\]
が全ての $n$ と $x$ で成り立つ。以下、$p$ と1の大小で場合を分ける。

$p>1$ のとき。式 ($**$) を積分し、仮定 $\int_Xf^p\,\d\mu=1$ を使うと
\[
0\le\int_X n\left|\sin\left(\frac{f(x)}n\right)\right|^p\,\d\mu
\le n^{1-p}\int_Xf^p\,\d\mu=n^{1-p}.
\]
$1-p<0$ だから $n^{1-p}\to0$ であり、はさみうちにより極限は $0$ である。

$p=1$ のとき。式 ($**$) は $n|\sin(f(x)/n)|\le f(x)$ となり、
$\int_Xf\,\d\mu=1<\infty$ だから $f$ は可積分な優関数である。
一方 式 ($*$) で $n^{1-p}=1$ なので、被積分関数は各点で $f(x)$ へ収束する。
よって優収束定理を適用でき
\[
\lim_{n\to\infty}\int_X n\left|\sin\left(\frac{f(x)}n\right)\right|\,\d\mu
=\int_Xf\,\d\mu=1.
\]

$0<p<1$ のとき。$1-p>0$ だから $n^{1-p}\to\infty$ であり、$f(x)^p>0$ と
式 ($*$) から、被積分関数は各点で $\infty$ へ発散する。
被積分関数は非負なので Fatou の補題が使えて
\[
\liminf_{n\to\infty}\int_X n\left|\sin\left(\frac{f(x)}n\right)\right|^p\,\d\mu
\ge\int_X\liminf_{n\to\infty}
n\left|\sin\left(\frac{f(x)}n\right)\right|^p\,\d\mu
=\int_X\infty\,\d\mu.
\]
ここで $\int_Xf^p\,\d\mu=1\ne0$ より $\mu(X)>0$ だから右辺は $\infty$ である。
従って極限は $+\infty$ である。

以上をまとめると、求める極限は
\[
\lim_{n\to\infty}\int_X n\left|\sin\left(\frac{f(x)}n\right)\right|^p\,\d\mu
=\begin{cases}
+\infty,&0<p<1,\\
1,&p=1,\\
0,&p>1
\end{cases}
\]
である。
\end{proof}

\begin{problem}{第8問｜独立確率変数の級数}{2015-s-8}
実数値確率変数 $X$ は連続型で，確率密度関数を $f$ とする。さらに，$X$ は2乗可積分であるとし，
$X$ の平均 $E[X]$ を $\mu$，分散 $V[X]$ を $\sigma^2$ とおく。
実数値確率変数列 $\{X_n\}_{n=1}^\infty$ は独立で，各 $X_n$ の分布関数は
\[
\int_{-\infty}^xnf(ny)\,\d y\qquad(x\in\R)
\]
によって与えられるものとする。整数 $n\ge1$ に対して
\[
S_n=\sum_{k=1}^nX_k
\]
とおくとき，以下の問いに答えよ。
\begin{enumerate}
  \item[(1)] 整数 $n\ge1$ に対して，$X_n$ の平均 $E[X_n]$ および分散 $V[X_n]$ を求めよ。
  \item[(2)] 確率変数列 $\{S_n\}_{n=1}^\infty$ が概収束するための必要十分条件は $\mu=0$ であることを示せ。
  \item[(3)] $X$ が標準正規分布にしたがうとき，(2) より $\{S_n\}_{n=1}^\infty$ はある確率変数 $S$ に概収束することがわかる。$S$ のしたがう確率分布を求めよ。必要があれば，
  \[
  \sum_{n=1}^\infty\frac1{n^2}=\frac{\pi^2}6
  \]
  を用いてもよい。
\end{enumerate}
\end{problem}
\begin{memo*}
\textbf{解答の見通し。}
密度の変数変換から $X_n$ は $X/n$ と同分布だと読み取る。
平均を引いた $Y_n=X_n-\mu/n$ は分散の総和が有限なので、その級数は概収束する。
残る平均部分は調和級数であり、$\mu=0$ かどうかを判定する。
標準正規の場合は独立な正規分布の和なので、分散の極限から極限分布が分かる。
\end{memo*}
\begin{proof}[解答]
(1) $X_n$ の分布関数を変数変換すると、任意の有界可測関数 $\psi$ に対し
\[
E[\psi(X_n)]=\int_{\R}\psi(x)nf(nx)\,\d x
=\int_{\R}\psi(u/n)f(u)\,\d u.
\]
従って $X_n$ は $X/n$ と同分布であり
\[
E[X_n]=\frac\mu n,\qquad V[X_n]=\frac{\sigma^2}{n^2}.
\]

(2) $Y_n=X_n-\mu/n$ とおくと、$Y_n$ は独立、平均0で
\[
\sum_{n=1}^\infty V[Y_n]
=\sigma^2\sum_{n=1}^\infty\frac1{n^2}<\infty.
\]
Kolmogorov の収束定理により $\sum_nY_n$ はほとんど確実に収束する。
従って
\[
S_n=\sum_{k=1}^nY_k+\mu\sum_{k=1}^n\frac1k.
\]
$\mu=0$ なら右辺第1項だけなので $S_n$ は概収束する。
$\mu\ne0$ なら第1項は有限値へ収束する一方、調和級数を含む第2項は
$\operatorname{sgn}(\mu)\infty$ へ発散するため、$S_n$ は有限な確率変数へ概収束しない。

(3) $X$ が標準正規分布なら $X_n\sim N(0,1/n^2)$ であり、独立性から
\[
S_n\sim N\left(0,\sum_{k=1}^n\frac1{k^2}\right).
\]
分散は $\sum_{k=1}^n1/k^2\to\pi^2/6$ だから、概収束極限の分布は
$N(0,\pi^2/6)$ である。
\end{proof}

\section{第2期・専門基礎科目（3題必修）}

\begin{problem}{第1問｜像・核と直和}{2015-2-1}
\begin{enumerate}
  \item[(1)] 
  \[
  A=\begin{pmatrix}1&2&3&0\\-1&0&-1&2\\2&3&5&-1\\1&1&2&-1\end{pmatrix},\quad
  B=\begin{pmatrix}1&2&3&0\\1&2&2&1\\2&4&3&3\\3&6&7&2\end{pmatrix}
  \]
  とし，列ベクトル空間 $\R^4$ のベクトル $v$ に対して $f_A(v)=Av$, $f_B(v)=Bv$ とおく。\\
  (i) $\operatorname{Im}f_A$ および $\operatorname{Im}f_B$ の基底を1組ずつ求めよ。また，$\operatorname{Im}f_A\cap\operatorname{Im}f_B$ を求めよ。\\
  (ii) $\dim(\ker f_A+\ker f_B)$ を求めよ。
  \item[(2)] $n$ を自然数とする。$f,g$ を $\R^n$ から $\R^n$ への線形写像とし，$\R^n=\operatorname{Im}f\oplus\operatorname{Im}g$ とする。\\
  (i) $\ker(f+g)=\ker f\cap\ker g$ であることを示せ。\\
  (ii) $\R^n=\ker f+\ker g$ であるとき，$f+g$ は全単射であることを示せ。
\end{enumerate}
\end{problem}
\begin{memo*}
\textbf{解答の見通し。}
行列の像は列ベクトルの全ての線形結合、核は $Ax=0$ を満たすベクトル全体である。
まず各列を少数の独立な列の線形結合として式で表し、その独立な列を像の基底にする。
共通部分のベクトルは $A$ 側と $B$ 側の二通りに基底表示できるため、二表示を等置して係数を解く。
核については、列の線形関係を $Ax$ の式へ代入して連立方程式を作る。

$U\oplus W$ は $U+W$ に加えて $U\cap W=\{0\}$ を満たす直和を意味する。
従って $f(x)=-g(x)$ が両方の像に属すれば0でなければならない。
最後は部分空間の和の次元公式を書き下し、$\ker(f+g)=\{0\}$、すなわち単射を示す。
同じ有限次元の空間同士では単射と全射が同値なので、全単射が従う。
\end{memo*}
\begin{proof}[解答]
\begin{enumerate}
  \item[(1)] (i) $A,B$ の列を調べると
  \[
  \operatorname{Im}f_A
  =\left\langle
  \begin{pmatrix}1\\-1\\2\\1\end{pmatrix},
  \begin{pmatrix}2\\0\\3\\1\end{pmatrix}
  \right\rangle,\qquad
  \operatorname{Im}f_B
  =\left\langle
  \begin{pmatrix}1\\1\\2\\3\end{pmatrix},
  \begin{pmatrix}3\\2\\3\\7\end{pmatrix}
  \right\rangle.
  \]
  これら4ベクトルの係数を成分比較すると、両方の像に属するベクトルの係数は
  全て0となる。従って
  \[
  \operatorname{Im}f_A\cap\operatorname{Im}f_B=\{0\}.
  \]

  (ii) 同じ列の関係から
  \[
  \ker f_A=\langle(-1,-1,1,0),(2,-1,0,1)\rangle,
  \]
  \[
  \ker f_B=\langle(-2,1,0,0),(-3,0,1,1)\rangle.
  \]
  $p=(-1,-1,1,0)$, $q=(2,-1,0,1)$, $r=(-2,1,0,0)$ とすると、
  第4のベクトルは $p+q+2r$ であり、$p,q,r$ は一次独立である。従って
  \[
  \dim(\ker f_A+\ker f_B)=3.
  \]

  \item[(2)] (i) $x\in\ker(f+g)$ なら
  $f(x)=-g(x)\in\operatorname{Im}f\cap\operatorname{Im}g=\{0\}$。
  従って $x\in\ker f\cap\ker g$。逆の包含は明らかだから
  \[
  \ker(f+g)=\ker f\cap\ker g.
  \]

  (ii) 階数・退化次数の定理と像の直和から
  \[
  n=\dim\operatorname{Im}f+\dim\operatorname{Im}g
  =2n-\dim\ker f-\dim\ker g.
  \]
  一方、$\R^n=\ker f+\ker g$ だから
  \[
  n=\dim\ker f+\dim\ker g-\dim(\ker f\cap\ker g).
  \]
  よって $\ker f\cap\ker g=\{0\}$。(i) より $f+g$ は単射であり、
  同じ有限次元空間の間の線形写像なので全単射である。
\end{enumerate}
\end{proof}

\begin{problem}{第2問｜滑らかな関数の分解}{2015-2-2}
$f(x,y)$ を $\R^2$ における $C^\infty$ 級関数で $f(0,0)=0$ を満たすものとする。このとき，次の問いに答えよ。
\begin{enumerate}
  \item[(1)] $f(x,y)=\displaystyle\int_0^1\frac{\d}{\d t}f(tx,ty)\,\d t$ を用いて
  \[
  f(x,y)=xg(x,y)+yh(x,y)
  \]
  となる $\R^2$ における $C^\infty$ 級関数 $g(x,y),h(x,y)$ が存在することを示せ。
  \item[(2)] $\R^2$ における $C^\infty$ 級関数 $G(x,y)$ と $\R$ における $C^\infty$ 級関数 $H(y)$ が存在して
  \[
  f(x,y)=xG(x,y)+yH(y)
  \]
  が成り立つことを示せ。
\end{enumerate}
\end{problem}
\begin{memo*}
\textbf{解答の見通し。}
目標の $f=xg+yh$ では係数 $x,y$ を作る必要がある。
そこで原点から $(x,y)$ へ向かう線分 $(tx,ty)$ を考え、
$\varphi(t)=f(tx,ty)$ を連鎖律で微分する。すると
$\varphi'(t)=x\frac{\partial f}{\partial x}(tx,ty)+y\frac{\partial f}{\partial y}(tx,ty)$ となり、必要な $x,y$ が現れる。
$t=0$ から1まで積分すれば、微積分学の基本定理により $f(x,y)-f(0,0)$ が戻る。

二つ目では $H$ を $y$ だけの関数にする必要がある。そのため
$f(x,y)-f(0,y)$ を $x$ 方向だけに積分し、残る $f(0,y)-f(0,0)$ を
$y$ 方向だけに積分する。
\end{memo*}
\begin{proof}[解答]
\begin{enumerate}
  \item[(1)] 連鎖律と $f(0,0)=0$ より
  \begin{align*}
  f(x,y)
  &=\int_0^1\frac{\d}{\d t}f(tx,ty)\,\d t\\
  &=x\int_0^1f_x(tx,ty)\,\d t
   +y\int_0^1f_y(tx,ty)\,\d t.
  \end{align*}
  従って
  \[
  g(x,y)=\int_0^1f_x(tx,ty)\,\d t,\qquad
  h(x,y)=\int_0^1f_y(tx,ty)\,\d t
  \]
  とすればよい。積分記号下で任意回微分できるので $g,h\in C^\infty$。

  \item[(2)]
  \[
  f(x,y)-f(0,y)=x\int_0^1f_x(tx,y)\,\d t,\qquad
  f(0,y)=y\int_0^1f_y(0,ty)\,\d t.
  \]
  よって
  \[
  G(x,y)=\int_0^1f_x(tx,y)\,\d t,\qquad
  H(y)=\int_0^1f_y(0,ty)\,\d t
  \]
  とおけば $f=xG+yH$ であり、$G,H$ はともに $C^\infty$ 級である。
\end{enumerate}
\end{proof}

\begin{problem}{第3問｜コンパクト性と連結性}{2015-2-3}
\begin{enumerate}
  \item[(1)] $S$ をコンパクトな位相空間とし，$T$ を Hausdorff 空間とする。$f:S\to T$ を $S$ から $T$ への連続写像とする。$f(S)$ の1点 $p$ に対し，$f^{-1}(\{p\})$ はコンパクト集合であることを示せ。
  \item[(2)] $\R$ を全ての実数からなる集合とし，$\Q$ を全ての有理数からなる集合とする。$\R^2$ の部分集合 $X,Y$ を $X=\Q^2$, $Y=\R^2\setminus X$ で定める。$X,Y$ に $\R^2$ からの相対位相を導入するとき，$X,Y$ が連結かどうかを判定せよ。
\end{enumerate}
\end{problem}
\begin{memo*}
\textbf{解答の見通し。}
ファイバー $f^{-1}(\{p\})$ は、$p$ へ写る点を全て集めた集合である。
Hausdorff 空間では一点集合が閉、連続写像では閉集合の逆像が閉、
コンパクト空間の閉部分集合はコンパクト、という三段階を順に使う。

$\Q^2$ の非連結性は、$x$ 座標がある無理数 $c$ より小さい側と大きい側に分けて示す。
$c$ 自身は有理点に現れないため二集合で全体を覆える。
$Y=\R^2\setminus\Q^2$ では、各点は少なくとも一方の座標が無理数である。
その無理数座標を固定した縦線または横線は $Y$ の外へ出ないので、折れ線で共通点へ結べる。
\end{memo*}
\begin{proof}[解答]
\begin{enumerate}
  \item[(1)] Hausdorff 空間では $\{p\}$ は閉である。従って連続性より
  $f^{-1}(\{p\})$ はコンパクト空間 $S$ の閉部分集合であり、コンパクトである。

  \item[(2)] $X=\Q^2$ は非連結である。例えば
  \[
  U=X\cap\{x<\sqrt2\},\qquad V=X\cap\{x>\sqrt2\}
  \]
  は非自明な開分離を与える。

  一方 $Y=\R^2\setminus\Q^2$ は道連結である。無理数 $\alpha,\beta$ を固定する。
  $(x,y)\in Y$ について、$x$ が無理数なら
  \[
  (x,y)\longrightarrow(x,\beta)\longrightarrow(\alpha,\beta),
  \]
  $x$ が有理数なら $y$ は無理数なので
  \[
  (x,y)\longrightarrow(\alpha,y)\longrightarrow(\alpha,\beta)
  \]
  という折れ線が全て $Y$ 内にある。従って $Y$ は連結である。
\end{enumerate}
\end{proof}

\section{応用数理コース・専門基礎科目（4題必修）}

\begin{problem}{第1問｜連立方程式と行列}{2015-a-1}
次の問いに答えよ。
\begin{enumerate}
  \item[(1)] 次の連立1次方程式が解をもつときの定数 $a$ の値を求め，そのときのすべての解を求めよ。
  \[
  \begin{cases}
  x+2y+3z=1\\
  4x+5y+6z=-2\\
  7x+8y+9z=a
  \end{cases}
  \]
  \item[(2)] 3つの行列
  \[
  A=\begin{pmatrix}1&2&3\\2&3&1\\3&1&2\end{pmatrix},\quad
  B=\begin{pmatrix}1&2&3\\8&5&5\\3&1&2\end{pmatrix},\quad
  C=\begin{pmatrix}2&1&3\\3&2&1\\1&3&2\end{pmatrix}
  \]
  に対して，$XA=B$, $AY=C$ をみたす行列 $X,Y$ を求めよ。
  \item[(3)] $A$ を $n$ 次正方行列，$I$ を $n$ 次単位行列とするとき，次を証明せよ。\\
  (i) $A^2+A-2I=O$ ならば，$A$ は正則である。ただし $O$ は零行列とする。\\
  (ii) $A$ がべき零行列ならば，$A$ は正則でなく，$I+A$ と $I-A$ は正則である。
\end{enumerate}
\end{problem}
\begin{memo*}
\textbf{解答の見通し。}
連立方程式では、第2式から第1式の4倍、第3式から第1式の7倍を引いて $x$ を消去する。
残った二式の左辺は定数倍の関係になるため、右辺も同じ倍率でなければ解が存在しない。

行列の積は一般に順序を交換できない。$XA=B$ から $X$ を孤立させるには右から
$A^{-1}$、$AY=C$ から $Y$ を孤立させるには左から $A^{-1}$ を掛ける。
その前に $|A|\ne0$ を確認し、余因子行列から逆行列を計算する。
最後の正則性では、逆行列の候補を多項式恒等式から直接作り、左右から掛けて単位行列になることを確認する。
\end{memo*}
\begin{proof}[解答]
\begin{enumerate}
  \item[(1)] 第2式から第1式の4倍、第3式から第1式の7倍を引くと
  \[
  -3y-6z=-6,\qquad -6y-12z=a-7.
  \]
  従って解をもつのは $a=-5$ のときであり、全解は
  \[
  (x,y,z)=(t-3,\,2-2t,\,t)\qquad(t\in\R).
  \]

  \item[(2)]
  \[
  A^{-1}=\frac1{18}
  \begin{pmatrix}-5&1&7\\1&7&-5\\7&-5&1\end{pmatrix}
  \]
  だから
  \[
  X=BA^{-1}=\begin{pmatrix}1&0&0\\0&1&2\\0&0&1\end{pmatrix},
  \qquad
  Y=A^{-1}C=\begin{pmatrix}0&1&0\\1&0&0\\0&0&1\end{pmatrix}.
  \]

  \item[(3)] (i) $A^2+A=2I$ より
  \[
  A\frac{A+I}{2}=\frac{A+I}{2}A=I.
  \]
  従って $A^{-1}=(A+I)/2$。

  (ii) $A^m=O$ とする。$A$ が正則なら $A^m$ も正則となるので矛盾。
  また
  \[
  (I-A)(I+A+\cdots+A^{m-1})=I-A^m=I,
  \]
  \[
  (I+A)(I-A+\cdots+(-1)^{m-1}A^{m-1})=I-(-A)^m=I.
  \]
  従って $I-A$, $I+A$ は正則である。
\end{enumerate}
\end{proof}

\begin{problem}{第2問｜対角化・漸化式・直交行列}{2015-a-2}
\begin{enumerate}
  \item[(1)] $A=\begin{pmatrix}-2&1&2\\1&0&0\\0&1&0\end{pmatrix}$ とするとき，次の問いに答えよ。\\
  (i) $A$ の固有値と各固有値に対する固有空間を求めよ。\\
  (ii) $A$ が対角化可能かどうか理由を含めて答えよ。また，対角化可能ならば対角化し，対角化するための正則行列とその逆行列も求めよ。\\
  (iii) $A^n$ の1行目の成分をすべて求めよ。\\
  (iv) 漸化式
  \[
  x_n=-2x_{n-1}+x_{n-2}+2x_{n-3}\qquad(n\ge4)
  \]
  について，初期条件が $x_1=1$, $x_2=0$, $x_3=0$ の場合と，初期条件が $x_1=0$, $x_2=-1$, $x_3=1$ の場合の一般項 $x_n$ をそれぞれ求めよ。
  \item[(2)] $A$ を $n$ 次正方行列とし，
  \begin{enumerate}
    \item[(ア)] $A$ は直交行列である
    \item[(イ)] $\R^n$ に属する任意のベクトル $x$ に対して，$\norm{Ax}=\norm{x}$ である
    \item[(ウ)] $\R^n$ に属する任意のベクトル $x,y$ に対して，$\inner{Ax}{Ay}=\inner{x}{y}$ である
  \end{enumerate}
  とする。ただし $\R^n$ に属するベクトル $a,b$ に対して，$\inner{a}{b}$ は $a$ と $b$ の内積を表し，$\norm{a}=\sqrt{\inner{a}{a}}$ とする。このとき，次を証明せよ。\\
  (i) (ア) が成立するならば，(イ) が成立する。\\
  (ii) (イ) が成立するならば，(ウ) が成立する。\\
  (iii) (ウ) が成立するならば，(ア) が成立する。
\end{enumerate}
\end{problem}
\begin{memo*}
\textbf{解答の見通し。}
まず $|\lambda I-A|=0$ を因数分解して固有値を求め、
各 $\lambda$ を $(A-\lambda I)x=0$ へ代入して固有ベクトルを得る。
固有ベクトルを列に並べた $P$ では $AP=PD$ が成り立つため、左から $P^{-1}$ を掛けて
$P^{-1}AP=D$、さらに $A^n=PD^nP^{-1}$ とできる。

漸化式へ試行解 $x_n=r^n$ を代入すると、行列と同じ三根 $1,-1,-2$ を持つ特性方程式が現れる。
根が相異なるので一般解を三つの指数列の線形結合として置き、初期値を代入して係数を解く。
直交行列の同値条件では、$\norm{x}^2=\inner{x}{x}$ と
$\norm{x+y}^2$ の展開を使い、「行列条件・長さ・内積」を順につなぐ。
\end{memo*}
\begin{proof}[解答]
\begin{enumerate}
  \item[(1)] (i)
  \[
  \det(\lambda I-A)=(\lambda-1)(\lambda+1)(\lambda+2),
  \]
  \[
  E_1=\langle(1,1,1)^{\mathsf T}\rangle,\quad
  E_{-1}=\langle(1,-1,1)^{\mathsf T}\rangle,\quad
  E_{-2}=\langle(4,-2,1)^{\mathsf T}\rangle.
  \]

  (ii) 固有値が相異なるので対角化可能である。
  \[
  P=\begin{pmatrix}1&1&4\\1&-1&-2\\1&1&1\end{pmatrix},\qquad
  P^{-1}=\frac16\begin{pmatrix}1&3&2\\-3&-3&6\\2&0&-2\end{pmatrix},
  \]
  \[
  P^{-1}AP=\diag(1,-1,-2).
  \]

  (iii) $A^n=P\diag(1,(-1)^n,(-2)^n)P^{-1}$ より、第1行は
  \[
  \left(
  \frac16-\frac{(-1)^n}{2}+\frac43(-2)^n,\quad
  \frac{1-(-1)^n}{2},\quad
  \frac13+(-1)^n-\frac43(-2)^n
  \right).
  \]

  (iv) 特性方程式は $(r-1)(r+1)(r+2)=0$ だから
  \[
  x_n=\alpha+\beta(-1)^n+\gamma(-2)^n.
  \]
  各初期条件を代入すると
  \[
  (x_1,x_2,x_3)=(1,0,0):
  \quad x_n=\frac13-(-1)^n+\frac16(-2)^n,
  \]
  \[
  (x_1,x_2,x_3)=(0,-1,1):
  \quad x_n=-\frac13-\frac16(-2)^n.
  \]

  \item[(2)] (i) $A^{\mathsf T}A=I$ なら
  \[
  \|Ax\|^2=x^{\mathsf T}A^{\mathsf T}Ax=\|x\|^2.
  \]

  (ii) 長さを保つなら、実内積の極化恒等式より
  \[
  \langle Ax,Ay\rangle
  =\frac{\|A(x+y)\|^2-\|Ax\|^2-\|Ay\|^2}{2}
  =\langle x,y\rangle.
  \]

  (iii) 内積を保つなら、標準基底 $e_i,e_j$ に対して
  \[
  (A^{\mathsf T}A)_{ij}
  =\langle Ae_i,Ae_j\rangle
  =\langle e_i,e_j\rangle=\delta_{ij}.
  \]
  従って $A^{\mathsf T}A=I$、すなわち $A$ は直交行列である。
\end{enumerate}
\end{proof}

\begin{problem}{第3問｜初等微積分}{2015-a-3}
次の問いに答えよ。
\begin{enumerate}
  \item[(1)] 次の関数の導関数を求めよ。\\
  (i) $\log|\tan x|$\\
  (ii) $x^{\sin x}$ $(x>0)$
  \item[(2)] (i) $x+\sqrt{x^2+1}=t$ のとき，$x$ を $t$ だけを用いて表わせ。\\
  (ii) 次の定積分の値を求めよ。
  \[
  \text{(a)}\ \int_0^1\frac{\d x}{\sqrt{x^2+1}}
  \qquad
  \text{(b)}\ \int_0^1\sqrt{x^2+1}\,\d x
  \]
\end{enumerate}
\end{problem}
\begin{memo*}
\textbf{解答の見通し。}
$\log|\tan x|$ は $\tan x\ne0$ の区間ごとに、合成関数の微分公式
$(\log|u|)'=u'/u$ を使う。$x^{\sin x}$ は実数値として $x>0$ で考え、
$x^{\sin x}=\exp(\sin x\log x)$ と書いて指数関数と積の微分へ直す。

置換 $t=x+\sqrt{x^2+1}$ では、共役な式を掛けると
\[
(\sqrt{x^2+1}+x)(\sqrt{x^2+1}-x)=1
\]
となるため、$t^{-1}=\sqrt{x^2+1}-x$ が得られる。
$t$ と $t^{-1}$ を加減して $x$ と平方根を表し、$\d x$ と積分端点も全て $t$ へ変換する。
\end{memo*}
\begin{proof}[解答]
\begin{enumerate}
  \item[(1)] (i)
  \[
  \frac{\d}{\d x}\log|\tan x|
  =\frac{\sec^2x}{\tan x}
  =\frac1{\sin x\cos x}.
  \]
  (ii) $x^{\sin x}=e^{\sin x\log x}$ より
  \[
  \frac{\d}{\d x}x^{\sin x}
  =x^{\sin x}\left(\cos x\log x+\frac{\sin x}{x}\right).
  \]

  \item[(2)] (i)
  \[
  t^{-1}=\sqrt{x^2+1}-x,\qquad
  x=\frac{t-t^{-1}}2.
  \]

  (ii) さらに
  \[
  \sqrt{x^2+1}=\frac{t+t^{-1}}2,\qquad
  \d x=\frac{1+t^{-2}}2\,\d t,
  \]
  で、$x=0,1$ は $t=1,1+\sqrt2$ に対応する。従って
  \[
  \text{(a)}\quad
  \int_0^1\frac{\d x}{\sqrt{x^2+1}}
  =\int_1^{1+\sqrt2}\frac{\d t}{t}
  =\log(1+\sqrt2),
  \]
  \[
  \text{(b)}\quad
  \int_0^1\sqrt{x^2+1}\,\d x
  =\frac14\int_1^{1+\sqrt2}(t+2t^{-1}+t^{-3})\,\d t
  =\frac{\sqrt2+\log(1+\sqrt2)}2.
  \]
\end{enumerate}
\end{proof}

\begin{problem}{第4問｜最大最小と三重積分}{2015-a-4}
次の問いに答えよ。
\begin{enumerate}
  \item[(1)] $\sqrt x+\sqrt y=1$ の条件の下で，$x^2+xy+y^2$ の最大値と最小値を求めよ。
  \item[(2)] 次の積分の値を求めよ。
  \[
  \iiint_V(x+y+z)\,\d x\d y\d z
  \]
  ただし $V=\{(x,y,z)\mid x^2+y^2+z^2\le a^2\}$ $(a>0)$ である。
\end{enumerate}
\end{problem}
\begin{memo*}
\textbf{解答の見通し。}
平方根が現れるので $u=\sqrt x$, $v=\sqrt y$ と置く。このとき $x=u^2$, $y=v^2$、
$u,v\ge0$ で、制約は直線 $u+v=1$ になる。
目的関数を展開し、$u+v=1$ を使って積 $q=uv$ だけの式へ直す。
相加相乗平均から $0\le q\le1/4$ なので、一変数二次関数の増減として最大・最小を決める。

三重積分では、球内の各点 $(x,y,z)$ と反対側の点 $(-x,-y,-z)$ が対になり、
領域の体積要素は変わらない一方、$x+y+z$ の符号だけが反転する。
変数変換を式にすると積分値 $J$ が $J=-J$ を満たし、$J=0$ となる。
\end{memo*}
\begin{proof}[解答]
\begin{enumerate}
  \item[(1)] $u=\sqrt x$, $v=\sqrt y$, $q=uv$ とおくと
  $u,v\ge0$, $u+v=1$, $0\le q\le1/4$ であり、
  \[
  x^2+xy+y^2=u^4+u^2v^2+v^4=1-4q+3q^2.
  \]
  これは $0\le q\le1/4$ で単調減少する。従って
  \[
  \max=1\quad ((x,y)=(1,0),(0,1)),
  \]
  \[
  \min=\frac3{16}\quad ((x,y)=(1/4,1/4)).
  \]

  \item[(2)] $V$ は原点対称で、$x+y+z$ は奇関数だから
  \[
  \iiint_V(x+y+z)\,\d x\d y\d z=0.
  \]
\end{enumerate}
\end{proof}
